\chapter{The different and trace estimates}

Throughout this chapter $E'/E$ is a finite separable extension of
complete discretely valued fields, with actual integer rings $B/A$.
The residue fields need not be finite.

\begin{definition}[The different and codifferent]
\label{fw:different}
\source{serre-1979-local-fields:page-0057}
\lean{differentIdeal, coeIdeal_differentIdeal}
\mathlibok
The codifferent is the fractional $B$-ideal
\[
 \mathfrak D_{B/A}^{-1}
 =\{x\in E':\operatorname{Tr}_{E'/E}(xB)\subseteq A\},
\]
and the different is its inverse, a nonzero integral ideal of $B$.
Write $d(E'/E)=v_{E'}(\mathfrak D_{B/A})$.
\end{definition}

\begin{lemma}[The conductor and the different]
\label{fw:conductor-different}
\uses{fw:different}
\lean{conductor_mul_differentIdeal}
\mathlibok
If $\theta\in B$ generates $E'$ over $E$, and $f\in A[X]$ is its
monic minimal polynomial, let
$\mathfrak f=\{b\in B:bB\subseteq A[\theta]\}$. Then
$\mathfrak f\mathfrak D_{B/A}=f'(\theta)B$.
\end{lemma}

\begin{lemma}[An integral generator generates the fraction field]
\label{fw:integral-fraction-generator}
\lean{FontaineWintenberger.adjoin_algebraMap_eq_top_of_integral_adjoin_eq_top}
\leanok
If $B=A[\theta]$, then $E'=E[\theta]$.
\end{lemma}
\begin{proof}
The field $E(\theta)$ contains the image of every element of
$A[\theta]=B$, and hence every quotient of two elements of $B$.
Since $E'$ is the fraction field of $B$, this proves
$E'=E(\theta)$. The element $\theta$ is algebraic over $E$,
so $E[\theta]$ is already a field and equals $E(\theta)$.
\end{proof}

\begin{lemma}[The different of a monogenic extension]
\label{fw:monogenic-different}
\source{serre-1979-local-fields:page-0063}
\dcref{Serre, Chapter III, Section 6, Corollary 2}
\uses{fw:conductor-different,fw:integral-fraction-generator,fw:totally-ramified-uniformizer-generation}
\lean{FontaineWintenberger.differentIdeal_eq_span_derivative_of_adjoin_eq_top, FontaineWintenberger.extensionValuationSubring_different_uniformizer}
\leanok
If $B=A[\theta]$ and $f\in A[X]$ is its monic minimal polynomial,
then $\mathfrak D_{B/A}=f'(\theta)B$.
In particular, in a totally ramified extension the formula applies to
every uniformizer of $B$.
\end{lemma}
\begin{proof}
The preceding lemma gives $E'=E[\theta]$. The conductor of
$A[\theta]=B$ in $B$ is the unit ideal, since every product of two
elements of $B$ is in $B$. The conductor--different identity
therefore gives $\mathfrak D_{B/A}=f'(\theta)B$.
Every uniformizer generates the integer ring in a totally ramified
extension, so the preceding argument applies to each such uniformizer.
\end{proof}

\begin{lemma}[Ideal orders and normalized integral valuations]
\label{fw:ideal-order-generator}
\uses{fw:ideal-power-valuation}
\lean{FontaineWintenberger.maximalIdeal_emultiplicity_span}
\leanok
For a discrete valuation ring $B$ and $x\in B$, the multiplicity
of its maximal ideal in $(x)$ is $v_B(x)$, including infinite
multiplicity for $x=0$.
\end{lemma}
\begin{proof}
For every integer $n\geq0$, $\mathfrak m_B^n$ divides $(x)$ exactly
when $(x)\subseteq\mathfrak m_B^n$, exactly when
$x\in\mathfrak m_B^n$, and hence exactly when $n\leq v_B(x)$.
These tests for all natural $n$ determine an element of
$\mathbb N\cup\{+\infty\}$ uniquely.
\end{proof}

\begin{lemma}[The different as a product of conjugate displacements]
\label{fw:different-conjugate-product}
\source{serre-1979-local-fields:page-0071}
\dcref{Serre, Chapter IV, Section 1, proof of Proposition 4}
\uses{fw:monogenic-different,fw:galois-generator-conjugates,fw:generator-evaluation,fw:ideal-order-generator}
\lean{FontaineWintenberger.aeval_derivative_minpoly_eq_prod_displacements, FontaineWintenberger.differentIdeal_eq_span_prod_displacements, FontaineWintenberger.differentIdeal_order_eq_sum_ramificationOrder}
\leanok
Suppose $E'/E$ is Galois with group $G$ and $B=A[\theta]$. Then
\[
 \mathfrak D_{B/A}
 =\left(\prod_{g\in G,\ g\ne1}(\theta-g\theta)\right),\qquad
 d(E'/E)=\sum_{g\ne1}j_{E'}(g).
\]
\end{lemma}
\begin{proof}
The minimal polynomial satisfies
$f(X)=\prod_{g\in G}(X-g\theta)$. Separate its identity factor:
$f(X)=(X-\theta)\prod_{g\ne1}(X-g\theta)$.
The product rule, evaluated at $\theta$, gives
$f'(\theta)=\prod_{g\ne1}(\theta-g\theta)$.
The monogenic different formula proves the ideal equality.
The order of a principal ideal equals the normalized valuation of
its generator. Multiplicativity of that valuation gives the sum
of $v_{E'}(\theta-g\theta)$, and generator evaluation identifies
each summand with $j_{E'}(g)$.
\end{proof}

\begin{theorem}[Hilbert's different formula]
\label{fw:hilbert-different}
\source{serre-1979-local-fields:page-0071}
\dcref{Serre, Chapter IV, Section 1, Proposition 4}
\uses{fw:different-conjugate-product,fw:integral-generator,fw:ramification-filtration,fw:integral-order-function,fw:local-inertia-size}
\lean{FontaineWintenberger.sum_ramificationOrder_eq_sum_lower_card, FontaineWintenberger.exists_sum_ramificationOrder_eq_sum_lower_card, FontaineWintenberger.spectralIntegerRing_different_order_eq_sum_ramificationOrder, FontaineWintenberger.spectralIntegerRing_different_order_eq_sum_lower_card}
\leanok
For a finite Galois extension $E'/E$ with group $G$ and perfect residue
field $k_E$,
\[
 d(E'/E)=\sum_{g\ne1}j_{E'}(g)
         =\sum_{n\geq0}(|G_n|-1).
\]
\end{theorem}
\begin{proof}
Choose an integral generator and apply the preceding lemma for the
first equality. For nonidentity $g$, $j_{E'}(g)$ is a finite natural
integer, and it equals the number of positive integers $n+1$ with
$n+1\leq j_{E'}(g)$. For fixed $n$, this inequality holds exactly
for $g\in G_n$. Counting the same finite set of pairs $(g,n)$ first
by $g$ and then by $n$ proves the second equality. Finiteness follows
from the eventual triviality of the lower groups.
\end{proof}

\begin{lemma}[The trace ideal of an integral ideal]
\label{fw:trace-ideal}
\source{serre-1979-local-fields:page-0057,serre-1979-local-fields:page-0058,serre-1979-local-fields:page-0090}
\dcref{Serre, Chapter III, Section 3, Proposition 7; Chapter V, Section 3, Lemma 4}
\uses{fw:different,fw:normalized-valuation-scaling,fw:ideal-order-generator}
\lean{FontaineWintenberger.trace_integral_iff_le_codifferent, FontaineWintenberger.trace_principal_iff_le_mul_codifferent, FontaineWintenberger.trace_principal_iff_mul_different_le, FontaineWintenberger.intTrace_mem_principal_iff_mul_different_le, FontaineWintenberger.intTraceIdeal, FontaineWintenberger.intTraceIdeal_le_iff, FontaineWintenberger.ideal_eq_pow_emultiplicity, FontaineWintenberger.intTraceIdeal_pow_le_pow_iff, FontaineWintenberger.intTraceIdeal_pow, FontaineWintenberger.extensionValuationSubring_intTraceIdeal_pow}
\leanok
Put $e=e(E'/E)$ and $d=d(E'/E)$. For every integer $n\geq0$,
\[
 \operatorname{Tr}_{E'/E}(\mathfrak m_B^n)
       =\mathfrak m_A^{\lfloor(d+n)/e\rfloor}.
\]
\end{lemma}
\begin{proof}
The trace image is the actual $A$-module image under the integral
trace, hence an ideal of $A$. For a fractional $B$-ideal $J$ and
nonzero principal fractional $A$-ideal $I$,
the codifferent definition gives
\[
 \operatorname{Tr}(J)\subseteq I
 \ \Longleftrightarrow\
 \operatorname{Tr}(I^{-1}J)\subseteq A
 \ \Longleftrightarrow\
 I^{-1}J\subseteq\mathfrak D^{-1}
 \ \Longleftrightarrow\
 J\subseteq I\mathfrak D^{-1}.
\]
The middle equivalence holds because $I^{-1}J$ is a $B$-module:
requiring its trace to lie in $A$ also requires the trace after
multiplication by every element of $B$ to lie in $A$.
Multiplying the last inclusion by the nonzero different gives the
equivalent condition $J\mathfrak D\subseteq IB$. Apply this with
$J=\mathfrak m_B^n$ and $I=\mathfrak m_A^s$, for $s\geq0$.
Every nonzero ideal of a discrete valuation ring is a power of its
maximal ideal, with exponent its maximal-ideal multiplicity. The
condition is therefore
\[
 \mathfrak m_B^{n+d}\subseteq\mathfrak m_B^{es},
 \quad\text{equivalently}\quad es\leq d+n.
\]
Here $e>0$. If the trace ideal were zero, it would be contained in
every $\mathfrak m_A^s$, forcing this inequality for arbitrarily
large $s$, a contradiction. Write the trace ideal as
$\mathfrak m_A^t$. Its containment in $\mathfrak m_A^s$ is exactly
$s\leq t$, so the same criterion gives
$s\leq t\Longleftrightarrow es\leq d+n$ for every $s\geq0$.
Taking $s=t$ and $s=\lfloor(d+n)/e\rfloor$ proves
$t=\lfloor(d+n)/e\rfloor$.
\end{proof}
